\documentclass[10pt,a4paper]{amsart}
\usepackage{amsmath,amssymb,amsthm,graphicx}
\usepackage{verbatim}
\newenvironment{reference}{\comment}{\endcomment}
%\newenvironment{reference}{}{}
\newenvironment{slogan}{\comment}{\endcomment}
\newenvironment{history}{\comment}{\endcomment}

% For commutative diagrams we use Xy-pic
\usepackage[all]{xy}

% We use 2cell for 2-commutative diagrams.
\xyoption{2cell}
% Theorem environments.
%
\theoremstyle{plain}
\newtheorem{theorem}[subsection]{Theorem}
\newtheorem{proposition}[subsection]{Proposition}
\newtheorem{lemma}[subsection]{Lemma}

\theoremstyle{definition}
\newtheorem{definition}[subsection]{Definition}
\newtheorem{example}[subsection]{Example}
\newtheorem{exercise}[subsection]{Exercise}
\newtheorem{situation}[subsection]{Situation}

\theoremstyle{remark}
\newtheorem{remark}[subsection]{Remark}
\newtheorem{remarks}[subsection]{Remarks}

\numberwithin{equation}{subsection}

% Macros
%
\def\lim{\mathop{\mathrm{lim}}\nolimits}
\def\colim{\mathop{\mathrm{colim}}\nolimits}
\def\Spec{\mathop{\mathrm{Spec}}}
\def\Hom{\mathop{\mathrm{Hom}}\nolimits}
\def\Ext{\mathop{\mathrm{Ext}}\nolimits}
\def\SheafHom{\mathop{\mathcal{H}\!\mathit{om}}\nolimits}
\def\SheafExt{\mathop{\mathcal{E}\!\mathit{xt}}\nolimits}
\def\Sch{\mathit{Sch}}
\def\Mor{\mathop{\mathrm{Mor}}\nolimits}
\def\Ob{\mathop{\mathrm{Ob}}\nolimits}
\def\Sh{\mathop{\mathit{Sh}}\nolimits}
\def\NL{\mathop{N\!L}\nolimits}
\def\CH{\mathop{\mathrm{CH}}\nolimits}
\def\proetale{{pro\text{-}\acute{e}tale}}
\def\etale{{\acute{e}tale}}
\def\QCoh{\mathit{QCoh}}
\def\Ker{\mathop{\mathrm{Ker}}}
\def\Im{\mathop{\mathrm{Im}}}
\def\Coker{\mathop{\mathrm{Coker}}}
\def\Coim{\mathop{\mathrm{Coim}}}

% Boxtimes
%
\DeclareMathSymbol{\boxtimes}{\mathbin}{AMSa}{"02}

%
% Macros for moduli stacks/spaces
%
\def\QCohstack{\mathcal{QC}\!\mathit{oh}}
\def\Cohstack{\mathcal{C}\!\mathit{oh}}
\def\Spacesstack{\mathcal{S}\!\mathit{paces}}
\def\Quotfunctor{\mathrm{Quot}}
\def\Hilbfunctor{\mathrm{Hilb}}
\def\Curvesstack{\mathcal{C}\!\mathit{urves}}
\def\Polarizedstack{\mathcal{P}\!\mathit{olarized}}
\def\Complexesstack{\mathcal{C}\!\mathit{omplexes}}
% \Pic is the operator that assigns to X its picard group, usage \Pic(X)
% \Picardstack_{X/B} denotes the Picard stack of X over B
% \Picardfunctor_{X/B} denotes the Picard functor of X over B
\def\Pic{\mathop{\mathrm{Pic}}\nolimits}
\def\Picardstack{\mathcal{P}\!\mathit{ic}}
\def\Picardfunctor{\mathrm{Pic}}
\def\Deformationcategory{\mathcal{D}\!\mathit{ef}}
\usepackage{cleveref}
\crefname{originaldefinition}{original definition}{original definitions}
\Crefname{originaldefinition}{Original definition}{Original definitions}
\makeatletter
\newlabel{definition-sieve-generated}{{00Z1}{source}{}{sourcepointer-1}{}}
\newlabel{definition-sieve-generated@cref}{{[originaldefinition][0][]00Z1}{[1][0][]source}{}{}{}}
\makeatother
\crefname{originaldefinition}{original definition}{original definitions}
\Crefname{originaldefinition}{Original definition}{Original definitions}
\makeatletter
\newlabel{definition-sheaf-sets-topology}{{00Z8}{source}{}{sourcepointer-2}{}}
\newlabel{definition-sheaf-sets-topology@cref}{{[originaldefinition][0][]00Z8}{[1][0][]source}{}{}{}}
\makeatother
\crefname{originaldefinition}{original definition}{original definitions}
\Crefname{originaldefinition}{Original definition}{Original definitions}
\makeatletter
\newlabel{definition-sheaf-sets}{{00VM}{source}{}{sourcepointer-3}{}}
\newlabel{definition-sheaf-sets@cref}{{[originaldefinition][0][]00VM}{[1][0][]source}{}{}{}}
\makeatother
\crefname{originaldefinition}{original definition}{original definitions}
\Crefname{originaldefinition}{Original definition}{Original definitions}
\makeatletter
\newlabel{definition-site}{{00VH}{source}{}{sourcepointer-4}{}}
\newlabel{definition-site@cref}{{[originaldefinition][0][]00VH}{[1][0][]source}{}{}{}}
\makeatother
\crefname{originaldefinition}{original definition}{original definitions}
\Crefname{originaldefinition}{Original definition}{Original definitions}
\makeatletter
\newlabel{definition-topology}{{00Z4}{source}{}{sourcepointer-5}{}}
\newlabel{definition-topology@cref}{{[originaldefinition][0][]00Z4}{[1][0][]source}{}{}{}}
\makeatother
\title{The topology defined by the coverings of a site has the same sheaves}
\author{The Stacks Project Authors}
\date{stacks-project:sites.tex:local-snapshot-20261012; GNU FDL 1.2 or later}
\begin{document}
\maketitle
\noindent\textit{Editorial note.} All delivered mathematical text is the author's own, copied from the pinned Stacks Project chapter file sites-cohomology.tex; no statement or proof step was rewritten, reordered or removed.\par\medskip
\noindent\textit{Editorial note.} The proof is the author's complete argument; every step the author writes out is inside the counted span and no closing sentence is a placeholder.\par\medskip
\noindent\textit{Editorial note.} Every labelled reference occurring inside the counted spans is carried in assembly.external\_references as a non-load-bearing navigation pointer (same-chapter ones at their definition site, cross-chapter ones at the invocation inside the counted span), so no load-bearing mathematics is replaced by a stub. No bibliography citation occurs inside the counted spans.\par\medskip
\noindent\textit{Editorial note.} The checked chapter list for same-name collisions is the set of stacks chapters staged in this repository: chow, cohomology, divisors, etale-cohomology, flat, more-algebra, more-morphisms, morphisms, spaces-more-morphisms, varieties together with this card's chapter sites-cohomology.\par\medskip
\noindent\textit{Editorial note.} The preamble is the chapter's own retained minimal preamble (sources/stacks-sites-cohomology/retained-preamble.tex), a verbatim excerpt of the native preamble: the Xy-pic setup and the native theorem declarations with the mathematical macros.\par\medskip
\noindent\textit{Allowed standard prerequisites (editorial).}\begin{itemize}
\item The axioms of a site and of a topology. Same chapter, Definition definition-site and Definition definition-topology. Supplies the axioms matched against each other in the first part.
\item The sieve generated by a family of morphisms. Same chapter, Definition definition-sieve-generated. Supplies the sieve appearing in the definition of the topology J.
\item The sheaf conditions for a topology and for a site. Same chapter, Definition definition-sheaf-sets-topology and Definition definition-sheaf-sets. Supplies the two conditions compared in the second part.
\end{itemize}\medskip
\ifdefined\hypertarget\hypertarget{sourcepointer-1}{}\fi
\noindent\textit{Source reference.} Original definition 00Z1: sources/stacks-sites/pinned-original-source.tex, source lines 11468--11468. Navigation pointer invoked at proof line 11468; the counted span contains the pointer only and the cited result is not reproduced inside it, so the pointer is registered as a non-load-bearing navigation stub in the established form for Stacks material. No bibliography citation is involved.\par\medskip
\ifdefined\hypertarget\hypertarget{sourcepointer-2}{}\fi
\noindent\textit{Source reference.} Original definition 00Z8: sources/stacks-sites/pinned-original-source.tex, source lines 11472--11472. Navigation pointer invoked at proof line 11472; the counted span contains the pointer only and the cited result is not reproduced inside it, so the pointer is registered as a non-load-bearing navigation stub in the established form for Stacks material. No bibliography citation is involved.\par\medskip
\ifdefined\hypertarget\hypertarget{sourcepointer-3}{}\fi
\noindent\textit{Source reference.} Original definition 00VM: sources/stacks-sites/pinned-original-source.tex, source lines 11474--11474. Navigation pointer invoked at proof line 11474; the counted span contains the pointer only and the cited result is not reproduced inside it, so the pointer is registered as a non-load-bearing navigation stub in the established form for Stacks material. No bibliography citation is involved.\par\medskip
\ifdefined\hypertarget\hypertarget{sourcepointer-4}{}\fi
\noindent\textit{Source reference.} Original definition 00VH: sources/stacks-etale-cohomology/pinned-original-source.tex, source lines 526--526. Navigation pointer invoked at proof line 11481; the counted span contains the pointer only and the cited result is not reproduced inside it, so the pointer is registered as a non-load-bearing navigation stub in the established form for Stacks material. No bibliography citation is involved.\par\medskip
\ifdefined\hypertarget\hypertarget{sourcepointer-5}{}\fi
\noindent\textit{Source reference.} Original definition 00Z4: sources/stacks-sites/pinned-original-source.tex, source lines 11484--11484. Navigation pointer invoked at proof line 11484; the counted span contains the pointer only and the cited result is not reproduced inside it, so the pointer is registered as a non-load-bearing navigation stub in the established form for Stacks material. No bibliography citation is involved.\par\medskip
\begin{lemma}
\label{lemma-site-gives-topology}
Let $\mathcal{C}$ be a site with coverings $\text{Cov}(\mathcal{C})$.
For every object $U$ of $\mathcal{C}$, let $J(U)$ denote
the set of sieves $S$ on $U$ with the following property:
there exists a covering
$\{f_i : U_i \to U\}_{i\in I} \in \text{Cov}(\mathcal{C})$
so that the sieve $S'$ generated by the $f_i$ (see Definition
\ref{definition-sieve-generated}) is contained in $S$.
\begin{enumerate}
\item This $J$ is a topology on $\mathcal{C}$.
\item A presheaf $\mathcal{F}$ is a sheaf for this topology
(see Definition \ref{definition-sheaf-sets-topology})
if and only if it is a sheaf on the site (see
Definition \ref{definition-sheaf-sets}).
\end{enumerate}
\end{lemma}

\begin{proof}
To prove the first assertion we just note that axioms
(1), (2) and (3) of the definition of a site
(Definition \ref{definition-site})
directly imply the axioms
(3), (2) and (1) of the definition of a topology
(Definition \ref{definition-topology}). As an example we
prove $J$ has property (2). Namely, let $U$ be an object
of $\mathcal{C}$, let $S, S'$ be sieves on $U$ such that
$S \in J(U)$, and such that for every $V \to U$ in $S(V)$
we have $S' \times_U V \in J(V)$. By definition of $J(U)$
we can find a covering $\{f_i : U_i \to U\}$ of the site
such that $S$ the image of $h_{U_i} \to h_U$ is contained
in $S$. Since each $S'\times_U U_i$ is in $J(U_i)$ we
see that there are coverings $\{U_{ij} \to U_i\}$ of the
site such that $h_{U_{ij}} \to h_{U_i}$ is contained
in $S' \times_U U_i$. By definition of the base change
this means that $h_{U_{ij}} \to h_U$ is contained
in the subpresheaf $S' \subset h_U$. By axiom (2) for
sites we see that $\{U_{ij} \to U\}$ is a covering of
$U$ and we conclude that $S' \in J(U)$ by definition of $J$.

\medskip\noindent
Let $\mathcal{F}$ be a presheaf. Suppose that $\mathcal{F}$
is a sheaf in the topology $J$. We will show that $\mathcal{F}$
is a sheaf on the site as well. Let $\{f_i : U_i \to U\}_{i\in I}$
be a covering of the site. Let $s_i \in \mathcal{F}(U_i)$ be a
family of sections such that
$s_i|_{U_i \times_U U_j} = s_j|_{U_i \times_U U_j}$ for all
$i, j$. We have to show that there exists a unique section
$s \in \mathcal{F}(U)$ restricting back to the $s_i$ on the $U_i$.
Let $S \subset h_U$ be the sieve generated by the $f_i$.
Note that $S \in J(U)$ by definition. In stead of constructing
$s$, by the sheaf condition in the topology, it suffices
to construct an element
$$
\varphi \in \Mor_{\textit{PSh}(\mathcal{C})}(S, \mathcal{F}).
$$
Take $\alpha \in S(T)$ for some object $T \in \mathcal{U}$.
This means exactly that $\alpha : T \to U$ is a morphism
which factors through $f_i$ for some $i\in I$ (and maybe more than $1$).
Pick such an index $i$ and a factorization $\alpha = f_i \circ \alpha_i$.
Define $\varphi(\alpha) = \alpha_i^* s_i$. If $i'$, $\alpha
= f_i \circ \alpha_{i'}'$ is a second choice, then
$\alpha_i^* s_i = (\alpha_{i'}')^* s_{i'}$ exactly because of our
condition $s_i|_{U_i \times_U U_j} = s_j|_{U_i \times_U U_j}$ for all
$i, j$. Thus $\varphi(\alpha)$ is well defined. We leave it to the reader
to verify that $\varphi$, which in turn determines $s$ is correct
in the sense that $s$ restricts back to $s_i$.

\medskip\noindent
Let $\mathcal{F}$ be a presheaf. Suppose that $\mathcal{F}$
is a sheaf on the site $(\mathcal{C}, \text{Cov}(\mathcal{C}))$.
We will show that $\mathcal{F}$ is a sheaf for the topology $J$
as well. Let $U$ be an object of $\mathcal{C}$. Let
$S$ be a covering sieve on $U$ with respect to the topology $J$.
Let
$$
\varphi \in \Mor_{\textit{PSh}(\mathcal{C})}(S, \mathcal{F}).
$$
We have to show there is a unique element in
$\mathcal{F}(U) = \Mor_{\textit{PSh}(\mathcal{C})}(h_U, \mathcal{F})$
which restricts back to $\varphi$. By definition there exists
a covering $\{f_i : U_i \to U\}_{i\in I} \in \text{Cov}(\mathcal{C})$
such that $f_i : U_i \to U$ belongs to $S(U_i)$.  Hence
we can set $s_i = \varphi(f_i) \in \mathcal{F}(U_i)$.
Then it is a pleasant exercise to see that
$s_i|_{U_i \times_U U_j} = s_j|_{U_i \times_U U_j}$ for all
$i, j$. Thus we obtain the desired section $s$ by the sheaf
condition for $\mathcal{F}$ on the site
$(\mathcal{C}, \text{Cov}(\mathcal{C}))$.
Details left to the reader.
\end{proof}

\end{document}
